Computer vision has been widely applied across sports analytics to extract performance metrics, identify tactical patterns, and assist coaching. Many systems rely on object detection models such as YOLO, which are effective at identifying small, fast-moving objects under varying lighting and camera conditions. YOLO’s single-shot architecture enables real-time inference and has been adopted in numerous sports-tracking applications. Temporal association methods such as SORT~\cite{sort} complement such detectors by linking frame-by-frame predictions into coherent trajectories using simple motion models. Together, detection and tracking form the basis for many sports ball-tracking pipelines.

In tennis, TrackNet~\cite{tracknet} introduced a heatmap-based ball localization approach that predicts per-pixel probability maps. This architecture enables robust tracking under strong motion blur and rapid directional changes, outperforming traditional bounding-box detection for small-object motion. However, TrackNet is tailored to tennis broadcasts, where ball scale is larger, background clutter is lower, and camera perspectives are standardized. Volleyball practice footage presents different challenges: players frequently occlude the ball, camera positions vary significantly, and the ball may occupy only a few pixels. These factors limit the direct applicability of TrackNet to detailed serve analysis.

Volleyball-specific computer vision research has expanded in recent years. Plankelj and Mlakar~\cite{plankelj2024} proposed a system using U-Net for court segmentation and YOLOv8 for ball detection, combined with homography to create tactical visualizations. Their results show that convolutional models can detect court geometry and ball trajectories even when image quality varies. However, their system focuses on match-level ball movement rather than serve mechanics, and it assumes visible court lines for segmentation, which is not guaranteed in typical baseline practice recordings.

More advanced volleyball research explores monocular 3D trajectory reconstruction. Dong et al.~\cite{dong2021} and similar works fuse temporal features with physics-based constraints to infer ball depth and full 3D trajectories. These approaches support applications such as set-type recognition or action prediction, but typically require high-quality video, predictable motion patterns, or multiple visible bounces. Their modeling complexity makes them less suitable for training environments where teams rely on handheld or smartphone recordings.

Commercial systems such as BallTimeAI~\cite{balltime} also provide automated volleyball analysis using proprietary models. These systems detect players and the ball, classify actions, and generate match summaries. However, they emphasize match analytics rather than technical skill evaluation. Because they are closed-source and rely on undisclosed datasets, they cannot be easily extended or validated by researchers, and they do not provide serve-level landing metrics or consistency measurements.

Despite progress in general ball tracking, volleyball match analysis, and monocular trajectory estimation, serve-focused analysis remains largely unexplored. Existing systems either require complex multi-camera setups, address general ball motion rather than serve mechanics, or do not provide quantitative landing metrics. The present work addresses this gap by combining a fine-tuned YOLO detector, lightweight tracking, and court calibration to produce serve-level trajectory and landing estimates from a single baseline camera.
